\section{Lecture 21 (March 5, 2025, est.): Nondegeneracy and Frobenius injectivity}
\label{lec21}
\noindent\textit{Scribe: Guillermo Sanmarco}

Let $k$ be algebraically closed of characteristic $p>0$ and let
$\mathcal C$ be a symmetric fusion category. We continue to use
\emph{nondegenerate} for $\dim(\mathcal C)\ne0$.

\subsection{Nondegenerate categories have super fiber functors}
\begin{proposition}\label{lec21-nondegenerate-fiber}
If $\mathcal C$ is nondegenerate, then it admits a symmetric tensor functor
to $\Vect$ when $p=2$, and to $\SVec$ when $p>2$.
\end{proposition}
\begin{proof}
Use Theorem~\ref{lec20-lifting} to lift $\mathcal C$ symmetrically over
$W(k)$. Extend the generic fiber to an algebraic closure of
$K=\operatorname{Frac}W(k)$. By Deligne's theorem for symmetric fusion
categories in characteristic zero, the result is
$\Rep_{\overline K}(G,z)$ for a finite group $G$ and a central element
$z$ with $z^2=1$. Here $z$ acts as parity and the symmetry is the super
symmetry. Its global dimension is $|G|$: the superdimensions of simples
differ from their ordinary dimensions only by signs, which disappear on
squaring.

Global dimension of the integral lifting reduces to
$\dim(\mathcal C)$. Thus $|G|$ is not divisible by $p$.
The category $\Rep_k(G,z)$ is consequently semisimple and nondegenerate,
and its standard integral lifting has the same geometric generic fiber.
Faithfulness of lifting, \cite[Corollary 9.9]{ENO}, identifies
$\mathcal C$ with $\Rep_k(G,z)$. The use of the geometric generic fiber
avoids assuming that the fraction field $K$ is algebraically closed.
For $p>2$, take the forgetful functor to $\SVec$. If $p=2$, the group has
odd order, so $z=1$ and the forgetful functor lands in $\Vect$.
\end{proof}

\subsection{A dimension test for full faithfulness}
For a tensor functor $F:\mathcal C\to\mathcal D$ between fusion
categories, write $F(\mathcal C)$ for its \emph{fusion image}: the full
subcategory closed under direct sums and direct summands generated by the
objects $F(X)$. This need not be merely the list of objects literally
equal to $F(X)$.

\begin{lemma}\label{lec21-image-dimension}
One has
\[
 \operatorname{FPdim}(F(\mathcal C))
 \leq\operatorname{FPdim}(\mathcal C),
\]
with equality if and only if $F$ is fully faithful.
\end{lemma}
\begin{proof}
Replace the target by the fusion image $\mathcal D_0=F(\mathcal C)$;
then $F$ is dominant. Let $R$ be its right adjoint. The standard dimension
identity for a dominant tensor functor is
\[
 \operatorname{FPdim}(\mathcal C)
 =\operatorname{FPdim}(\mathcal D_0)
   \operatorname{FPdim}(R(\mathbf1)).
\]
For completeness, if $N_{ia}$ is the multiplicity of the simple $Y_a$ in
$F(X_i)$, the vector $\sum_i\operatorname{FPdim}(X_i)N_{ia}$ is the
image of the regular dimension vector of $\mathcal C$. It is therefore
an eigenvector for multiplication by $F(\bigoplus_iX_i)$.
Dominance means that this object contains every simple of $\mathcal D_0$;
its multiplication matrix has strictly positive entries, by rigidity.
Perron--Frobenius uniqueness makes the displayed vector a multiple of
$\operatorname{FPdim}(Y_a)$; its coefficient at the unit is
$\operatorname{FPdim}(R(\mathbf1))$. Summing against
$\operatorname{FPdim}(Y_a)$ gives the identity.
The unit occurs once in $R(\mathbf1)$ by adjunction, so this last
dimension is at least one. Equality holds exactly when $R(\mathbf1)$ is
the unit. The projection formula
$R(F(X)\otimes Y)\simeq X\otimes R(Y)$, obtained by rigidity and
adjunction, then gives $RF\simeq\mathrm{id}$, which is equivalent to
full faithfulness. Conversely a fully faithful $F$ gives
$R(\mathbf1)\simeq\mathbf1$.
\end{proof}

\begin{definition}
The category $\mathcal C$ is \emph{Frobenius injective} if
$\mathrm{Fr}:\mathcal C\to\mathcal C^{(1)}\boxtimes\mathsf{Ver}_p$ is
fully faithful.
\end{definition}
The word ``injective'' in this definition means fully faithful, not merely
faithful. Every tensor functor under consideration is already faithful.

\subsection{A Grothendieck-ring criterion}
\begin{lemma}\label{lec21-ring-criterion}
If the finite-dimensional algebra
$A=K_0(\mathcal C)\otimes_{\mathbb Z}k$ is semisimple, then
$\mathcal C$ is nondegenerate.
\end{lemma}
\begin{proof}
Let $t(a)$ be the trace of left multiplication by $a$ on $A$.
Since $A$ is commutative semisimple over the algebraically closed field
$k$, it is a product of copies of $k$, and $(a,b)\mapsto t(ab)$ is
nondegenerate. Set $r=\sum_Z[Z^*\otimes Z]$, summing over the simples.
For simple $X,Y$, rigidity gives
\[
 t([X][Y])
 =\sum_Z\dim_k\Hom(X\otimes Y\otimes Z,Z)
 =\dim_k\Hom(X\otimes Y,\bigoplus_Z Z^*\otimes Z),
\]
with the integer on the right read in $k$.
The matrix of multiplication by $r$ therefore agrees with the matrix of
this trace pairing after permuting the basis by duality. It is invertible,
so $r$ is a unit of $A$. Apply the categorical-dimension homomorphism
$A\to k$: the image of $r$ is
$\sum_Z\dim(Z)^2=\dim(\mathcal C)$, which must be nonzero.
\end{proof}

\begin{lemma}[Type $\Vect$]\label{lec21-type-vec}
A Frobenius-injective category of Frobenius type $\Vect$ is nondegenerate.
\end{lemma}
\begin{proof}
Here $\mathrm{Fr}:\mathcal C\to\mathcal C^{(1)}$ is fully faithful.
The two categories have the same Frobenius--Perron dimension, so this is
an equivalence. Equation~\eqref{lec20-type-vec-power} shows that taking
$p$th powers permutes the simple-object basis of $A$, with the
Frobenius automorphism on scalar coefficients. Thus the map
$a\mapsto a^p$ is bijective. It is Frobenius-semilinear, rather than
$k$-linear, but bijectivity still follows because $k$ is perfect.
An element killed by some power $p^m$ must therefore be zero. Hence
$A$ is reduced, and a finite-dimensional commutative reduced algebra
over $k$ is semisimple. Apply Lemma~\ref{lec21-ring-criterion}.
\end{proof}

\begin{lemma}[Type $\SVec$]\label{lec21-type-svec}
Suppose $p>2$. A Frobenius-injective category of Frobenius type $\SVec$
is nondegenerate.
\end{lemma}
\begin{proof}
Full faithfulness makes the image of a simple object simple. Thus for
every simple $X$, either
$\mathrm{Fr}(X)=Y\boxtimes\mathbf1$ or
$\mathrm{Fr}(X)=Y\boxtimes L_{p-1}$.
These alternatives give a faithful $\mathbb Z/2$-grading
$\mathcal C=\mathcal C_0\oplus\mathcal C_1$. The neutral component has
Frobenius type $\Vect$, and its Frobenius functor is fully faithful by
naturality. Lemma~\ref{lec21-type-vec} gives
$\dim(\mathcal C_0)\ne0$.

Here is the dimension calculation for the grading. Put
$r_i=\sum_{X\in\mathcal C_i}\dim(X)[X]$ in $K_0(\mathcal C)\otimes k$.
For a simple $Y\in\mathcal C_1$, the fusion rules and rigidity imply
$[Y]r_0=\dim(Y)r_1$; this follows by comparing the coefficient of each
simple, using $N_{Y,X}^{Z}=N_{Y^*,Z}^{X}$.
Apply categorical dimension and cancel $\dim(Y)\ne0$ to obtain
$\dim(\mathcal C_0)=\dim(\mathcal C_1)$. Consequently
\[
 \dim(\mathcal C)=2\dim(\mathcal C_0)\ne0.
\]
\end{proof}

\subsection{The other Frobenius types}
\begin{lemma}\label{lec21-ver-plus-in-image}
Suppose $p>3$ and $\mathcal C$ has Frobenius type $\mathsf{Ver}_p^+$ or
$\mathsf{Ver}_p$. Then its Frobenius fusion image contains
$\mathbf1\boxtimes\mathsf{Ver}_p^+$.
\end{lemma}
\begin{proof}
Some $\mathrm{Fr}(X)$ contains a simple summand $Y\boxtimes L_r$ with
$r\ne1,p-1$. Its tensor product with its dual contains
$(Y\otimes Y^*)\boxtimes(L_r\otimes L_r^*)$.
The first factor contains the unit, and the fusion rules show that the
second contains $L_3$. Thus the fusion image contains
$\mathbf1\boxtimes L_3$, which generates
$\mathbf1\boxtimes\mathsf{Ver}_p^+$.
\end{proof}
This lemma does not require Frobenius injectivity. Full faithfulness will
be used in the next lecture to lift individual summands back to
$\mathcal C$ and complete the proof of Ostrik's theorem.
